% Generálva: html/build_tex.js (ai01.js). Kézzel ne szerkeszd, a forrás a html/data mappa.
\setcounter{chapter}{0}
\chapter{\texorpdfstring{Jelek, teljesítmény, csúcstényező}{Jelek, teljesítmény, csúcstényező}}\label{f1}
Rövid, de a vizsgán gyakori kérdések: mekkora egy összetett jel csúcsértéke, teljesítménye és csúcstényezője. A kulcs: azonos frekvenciájú összetevők egyetlen szinusszá adódnak össze, különböző frekvenciájúak nem.

\section{Alapok, egyszerűen}
\subsection{\texorpdfstring{A1. Jelölések: teljesítmény és csúcstényező}{A1. Jelölések: teljesítmény és csúcstényező}}\label{a1-teljesitmeny}
\par\noindent{\small\renewcommand{\arraystretch}{1.25}\begin{tabularx}{\linewidth}{@{}>{\raggedright\arraybackslash}X>{\raggedright\arraybackslash}X@{}}\toprule
\textbf{Jel} & \textbf{Jelentés}\\\midrule
undefined & csúcsérték: a jel legnagyobb abszolút értéke\\
undefined & átlagteljesítmény 1 $\Omega$-on: a jel négyzetének időátlaga, undefined\\
undefined & effektív érték, undefined\\
undefined & csúcstényező, undefined\\
\bottomrule\end{tabularx}}\par

Egy $A$ amplitúdójú szinusz teljesítménye $A^2/2$, mert $\sin^2$ átlaga egy perióduson $1/2$. Ezért szinuszra $C=A/(A/\sqrt2)=\sqrt2$. Négyszögjelre $C=1$, háromszögjelre $\sqrt3$.

Különböző frekvenciájú szinuszok szorzatának átlaga 0, ezért a teljesítményük \textbf{összeadódik}: $P=\sum A_i^2/2$. Ugyanez igaz egy azonos frekvenciájú $\sin$ és $\cos$ párra is.

\subsection{\texorpdfstring{A2. Azonos frekvenciájú összetevők összeadása}{A2. Azonos frekvenciájú összetevők összeadása}}\label{a1-fazor}
Az addíciós tétellel $A\sin(\omega t+\varphi)=A\cos\varphi\sin\omega t+A\sin\varphi\cos\omega t$. Összevetve az $a\sin\omega t+b\cos\omega t$ alakkal: $A\cos\varphi=a$, $A\sin\varphi=b$, ezért

\begin{center}\adjustbox{max width=\linewidth}{$\displaystyle a\sin\omega t+b\cos\omega t=\sqrt{a^2+b^2}\,\sin(\omega t+\varphi),\qquad\mathrm{tg}\,\varphi=\frac ba$}\end{center}

Gondolj két egymásra merőleges nyílra: az eredő hossza Pitagorasz-tétellel jön ki. A csúcsérték ezért nem $a+b$, hanem $\sqrt{a^2+b^2}$.

Ha a frekvenciák aránya irracionális (pl. $\sqrt2$), a két szinusz előbb-utóbb tetszőlegesen közel kerül egyszerre a csúcsához, ezért a csúcsérték az amplitúdók \textbf{összege}.

% ------------------------------------------------------------------
\section{\texorpdfstring{1.1. feladat: $3\sin+4\cos$: csúcsérték, teljesítmény, csúcstényező}{1.1. feladat: 3sin+4cos: csúcsérték, teljesítmény, csúcstényező}}\label{ai1_01}
\setcounter{lepes}{0}
\begin{feladat}{1.1. feladat}
Határozza meg az $x(t)=3\sin(2\pi Ft)+4\cos(2\pi Ft)$ jel csúcstényezőjét, maximális értékét és teljesítményét!

{\footnotesize\sffamily Forrás: ALL-IN 90. oldal $\cdot$ aiD\_1\_1 (és 5-ösöm 2. oldal). Címkék: beugró, javítás}
\end{feladat}
\begin{keresunk}
\begin{itemize}
\item $C$: csúcstényező
\item $\hat x$: legnagyobb érték
\item $P$: átlagteljesítmény (1 $\Omega$-on)
\end{itemize}
\end{keresunk}
\begin{jeloles}
\begin{jeltabla}
$F$ & a jel frekvenciája & \\
$C$ & csúcstényező: csúcsérték osztva az effektív értékkel & \emph{keresett}\\
$\hat x$ & csúcsérték, a jel legnagyobb abszolút értéke & \emph{keresett}\\
$P$ & átlagteljesítmény 1 $\Omega$-on & \emph{keresett}\\
$a$ & a szinuszos tag együtthatója (3) & \\
$\omega$ & körfrekvencia, a frekvencia 2$\pi$-szerese & \\
$b$ & a koszinuszos tag együtthatója (4) & \\
$\varphi$ & az összevont szinusz kezdőfázisa & \\
$A$ & az összevont szinusz amplitúdója, 5 & \\
$x_\text{eff}$ & a jel effektív értéke & \\
\end{jeltabla}
\par\smallskip\textbf{Adott még:}
\begin{itemize}
\item $x(t)$: két azonos, $F$ frekvenciájú tag, a szinuszos amplitúdója 3, a koszinuszosé 4
\item $F$: a két tag közös frekvenciája. Az értéke nem kell, az eredmény nem függ tőle
\end{itemize}
\end{jeloles}
\begin{elmelet}[Miért így? Egyszerűen]
\textbf{Miért nem $3+4=7$ a csúcs?}\par Mert a szinusz és a koszinusz negyed periódussal el van tolva, soha nem érik el egyszerre a csúcsukat. Az eredő egyetlen szinusz (A2).
\end{elmelet}

\lepes{Egyetlen szinusszá alakítás}
Az $a\sin\omega t+b\cos\omega t=\sqrt{a^2+b^2}\,\sin(\omega t+\varphi)$ összevonással, ahol $a=3$, $b=4$ és $\omega=2\pi F$.

\begin{center}\adjustbox{max width=\linewidth}{$\displaystyle A=\sqrt{a^2+b^2}=\sqrt{3^2+4^2}=\sqrt{25}=5$}\end{center}
\jegyzet{tétel}{}

\begin{center}\adjustbox{max width=\linewidth}{$\displaystyle \mathrm{tg}\,\varphi=\tfrac43\ \Rightarrow\ \varphi=53{,}1^{\circ}$}\end{center}
\jegyzet{számérték}{$x(t)=5\sin(\omega t+53{,}1^{\circ})$}

\lepes{A maximális érték}
A szinusz legnagyobb értéke 1, ezért a jelé az amplitúdó.

\begin{center}\adjustbox{max width=\linewidth}{$\displaystyle \hat x=5$}\end{center}
\jegyzet{számérték}{}

\lepes{A teljesítmény}
A $\sin$ és a $\cos$ szorzatának átlaga 0, ezért a két tag teljesítménye összeadódik (A1). Ugyanez jön ki az összevont szinuszból is.

\begin{center}\adjustbox{max width=\linewidth}{$\displaystyle P=\frac{3^2}{2}+\frac{4^2}{2}=4{,}5+8=12{,}5\ \text{W}$}\end{center}
\jegyzet{definíció}{1 $\Omega$-on}

\begin{center}\adjustbox{max width=\linewidth}{$\displaystyle \frac{5^2}{2}=12{,}5$}\end{center}
\jegyzet{ellenőrzés}{egyezik}

\lepes{A csúcstényező}


\begin{center}\adjustbox{max width=\linewidth}{$\displaystyle x_\text{eff}=\sqrt{12{,}5}=3{,}54\qquad C=\frac{\hat x}{x_\text{eff}}=\frac{5}{3{,}54}=\sqrt2$}\end{center}
\jegyzet{definíció}{mint minden szinuszé}
\begin{tipp}[A gyűjtemény megoldása]
A csúcstényezőt ($\sqrt2$) és a teljesítményt jól adja, a maximumot viszont ott keresi, ahol $\sin=\cos$, és $7/\sqrt2\approx4{,}95$-öt kap.
\end{tipp}
\begin{figyelem}[Javítás]
A maximum nem a $\sin=\cos$ (45\textdegree{}-os) pontban van, hanem $\varphi=53{,}1^{\circ}$-nál. Helyesen $\hat x=5$.
\end{figyelem}
\begin{eredmeny}

\begin{center}\adjustbox{max width=\linewidth}{$\displaystyle C=\sqrt2\qquad\hat x=5\qquad P=12{,}5\ \text{W (1 }\Omega\text{-on)}$}\end{center}
\end{eredmeny}
\begin{figyelem}[Tipikus hibák]
\begin{itemize}
\item A két amplitúdó összeadása ($7$), vagy egy önkényesen választott időpontban vett érték.
\end{itemize}
\end{figyelem}


% ------------------------------------------------------------------
\section{\texorpdfstring{1.2. feladat: Két koszinusz, $f_1/f_2=\sqrt2$}{1.2. feladat: Két koszinusz, f1/f2=√2}}\label{ai1_02}
\setcounter{lepes}{0}
\begin{feladat}{1.2. feladat}
Mekkora az $x(t)=2\cos(2\pi f_1t)+2\cos(2\pi f_2t)$ jel csúcstényezője, legnagyobb értéke és teljesítménye, ha $f_1/f_2=\sqrt2$?

{\footnotesize\sffamily Forrás: ALL-IN 90. oldal $\cdot$ aiD\_1\_2. Címkék: beugró}
\end{feladat}
\begin{keresunk}
\begin{itemize}
\item $C$ (csúcstényező: csúcsérték osztva az effektív értékkel), $\hat x$ (csúcsérték, a jel legnagyobb abszolút értéke), $P$ (átlagteljesítmény 1 $\Omega$-on)
\end{itemize}
\end{keresunk}
\begin{jeloles}
\begin{jeltabla}
$f_1$ & az első koszinusz frekvenciája & \\
$f_2$ & a második koszinusz frekvenciája & \\
$C$ & csúcstényező: csúcsérték osztva az effektív értékkel & \emph{keresett}\\
$\hat x$ & csúcsérték, a jel legnagyobb abszolút értéke & \emph{keresett}\\
$P$ & átlagteljesítmény 1 $\Omega$-on & \emph{keresett}\\
$x$ & a vizsgált jel, x(t) & \\
\end{jeltabla}
\par\smallskip\textbf{Adott még:}
\begin{itemize}
\item két koszinusz, mindkettő 2 amplitúdójú
\end{itemize}
\end{jeloles}
\begin{elmelet}[Miért így? Egyszerűen]
\textbf{Miért $4$ a csúcs?}\par Az irracionális frekvenciaarány miatt a két fázis „mindenféle” párosítást bejár, köztük tetszőlegesen közel kerülnek a $(0,0)$ fázispárhoz, ahol mindkét koszinusz 1 (A2).
\end{elmelet}

\lepes{A legnagyobb érték}
Egyik tag sem lehet 2-nél nagyobb, és a két csúcs tetszőlegesen közel kerül egymáshoz.

\begin{center}\adjustbox{max width=\linewidth}{$\displaystyle \hat x=2+2=4$}\end{center}
\jegyzet{közelítés}{tetszőlegesen megközelíthető}

\lepes{A teljesítmény}
Különböző frekvenciák, a teljesítmények összeadódnak (A1).

\begin{center}\adjustbox{max width=\linewidth}{$\displaystyle P=\frac{2^2}{2}+\frac{2^2}{2}=2+2=4$}\end{center}
\jegyzet{definíció}{}

\lepes{A csúcstényező}


\begin{center}\adjustbox{max width=\linewidth}{$\displaystyle C=\frac{\hat x}{\sqrt P}=\frac{4}{\sqrt4}=2$}\end{center}
\jegyzet{definíció}{}
\begin{eredmeny}

\begin{center}\adjustbox{max width=\linewidth}{$\displaystyle \hat x=4\qquad P=4\ (\text{W, }1\,\Omega)\qquad C=2$}\end{center}
\end{eredmeny}
\begin{tipp}[Ellenőrzés]
\begin{itemize}
\item A gyűjtemény is 2, 4, 4 értéket ad. Ha a frekvenciák egyformák lennének, $x=4\cos$, $C=\sqrt2$ jönne ki: a különbség az összemérhetetlenségből adódik.
\end{itemize}
\end{tipp}


% ------------------------------------------------------------------
\section{\texorpdfstring{1.3. feladat: $4\sin-3\cos$: négy állítás}{1.3. feladat: 4sin-3cos: négy állítás}}\label{ai1_03}
\setcounter{lepes}{0}
\begin{feladat}{1.3. feladat}
Az $s(t)=4\sin(2\pi f_0t)-3\cos(2\pi f_0t)$ jel: A) csúcsértéke 4. B) csúcstényezője $1/\sqrt2$. C) $t=42/f_0$-ban az értéke $-4$. D) átlagteljesítménye 3,5. Melyik igaz?

{\footnotesize\sffamily Forrás: ALL-IN 91. oldal $\cdot$ aiD\_1\_3 (feleletválasztós 4.6.). Címkék: beugró, javítás}
\end{feladat}
\begin{keresunk}
\begin{itemize}
\item melyik állítás igaz
\end{itemize}
\end{keresunk}
\begin{jeloles}
\begin{jeltabla}
$s$ & a vizsgált jel, s(t) & \\
$f_0$ & a két tag közös frekvenciája & \\
$a$ & a szinuszos tag együtthatója (4) & \\
$\omega$ & körfrekvencia, a frekvencia 2$\pi$-szerese & \\
$b$ & a koszinuszos tag együtthatója ($-$3) & \\
$\varphi$ & az összevont szinusz kezdőfázisa & \\
$\hat s$ & a jel csúcsértéke & \\
$C$ & csúcstényező: csúcsérték osztva az effektív értékkel & \\
$P$ & átlagteljesítmény 1 $\Omega$-on & \\
\end{jeltabla}
\par\smallskip\textbf{Adott még:}
\begin{itemize}
\item azonos frekvenciájú szinusz (4) és koszinusz ($-3$)
\end{itemize}
\end{jeloles}

\lepes{A) csúcsérték}
Az $a\sin\omega t+b\cos\omega t=\sqrt{a^2+b^2}\,\sin(\omega t+\varphi)$ összevonással, ahol $a=4$, $b=-3$.

\begin{center}\adjustbox{max width=\linewidth}{$\displaystyle \hat s=\sqrt{4^2+(-3)^2}=5$}\end{center}
\jegyzet{tétel}{A hamis}

\lepes{B) csúcstényező}
Az összeg egyetlen szinusz, ezért $C=\sqrt2$. Az $1/\sqrt2$ már csak azért is rossz, mert $C\ge1$ mindig.

\begin{center}\adjustbox{max width=\linewidth}{$\displaystyle C=\frac{5}{5/\sqrt2}=\sqrt2$}\end{center}
\jegyzet{számérték}{B hamis}

\lepes{C) A jel értéke a $t=42/f_0$ időpontban}
Behelyettesítve az argumentum $2\pi\cdot42$, ez egész számú periódus.

\begin{center}\adjustbox{max width=\linewidth}{$\displaystyle s=4\sin(84\pi)-3\cos(84\pi)=4\cdot0-3\cdot1=-3$}\end{center}
\jegyzet{behelyettesítés}{C hamis}

\lepes{D) átlagteljesítmény}


\begin{center}\adjustbox{max width=\linewidth}{$\displaystyle P=\frac{4^2}{2}+\frac{3^2}{2}=8+4{,}5=12{,}5$}\end{center}
\jegyzet{definíció}{D hamis}
\begin{figyelem}[Javítás]
A gyűjtemény lapján A, B, C be van karikázva. Ha ez a „hamis” jelölése, a megoldás helyes, igaznak jelölve mindhárom rossz.
\end{figyelem}
\begin{eredmeny}

\begin{center}\adjustbox{max width=\linewidth}{$\displaystyle \text{egyik állítás sem igaz}$}\end{center}
\end{eredmeny}


% ------------------------------------------------------------------
\section{\texorpdfstring{1.4. feladat: $2\sin+2\cos$ csúcstényezője}{1.4. feladat: 2sin+2cos csúcstényezője}}\label{ai1_04}
\setcounter{lepes}{0}
\begin{feladat}{1.4. feladat}
Határozza meg az $x(t)=2\sin(2\pi ft)+2\cos(2\pi ft)$ jel csúcstényezőjét!

{\footnotesize\sffamily Forrás: ALL-IN 91. oldal $\cdot$ aiD\_1\_5. Címkék: beugró}
\end{feladat}
\begin{keresunk}
\begin{itemize}
\item $C$ (csúcstényező: csúcsérték osztva az effektív értékkel)
\end{itemize}
\end{keresunk}
\begin{jeloles}
\begin{jeltabla}
$C$ & csúcstényező: csúcsérték osztva az effektív értékkel & \emph{keresett}\\
$x$ & a vizsgált jel, x(t) & \\
$f$ & a két tag közös frekvenciája & \\
$a$ & a szinuszos tag együtthatója (2) & \\
$\omega$ & körfrekvencia, a frekvencia 2$\pi$-szerese & \\
$b$ & a koszinuszos tag együtthatója (2) & \\
$\varphi$ & az összevont szinusz kezdőfázisa & \\
$A$ & az összevont szinusz amplitúdója, $2\sqrt2$ & \\
$\hat x$ & csúcsérték, a jel legnagyobb abszolút értéke & \\
$P$ & átlagteljesítmény 1 $\Omega$-on & \\
\end{jeltabla}
\par\smallskip\textbf{Adott még:}
\begin{itemize}
\item azonos frekvenciájú szinusz és koszinusz, mindkettő 2 amplitúdójú
\end{itemize}
\end{jeloles}

\lepes{Egyetlen szinusz}
Az $a\sin\omega t+b\cos\omega t=\sqrt{a^2+b^2}\,\sin(\omega t+\varphi)$ összevonással, ahol $a=b=2$.

\begin{center}\adjustbox{max width=\linewidth}{$\displaystyle A=\sqrt{2^2+2^2}=2\sqrt2\qquad\mathrm{tg}\,\varphi=1\ \Rightarrow\ \varphi=45^{\circ}$}\end{center}
\jegyzet{tétel}{$x=2\sqrt2\sin(2\pi ft+45^{\circ})$}

\lepes{Csúcstényező}
Szinuszos jel, $\hat x=2\sqrt2$, $P=(2\sqrt2)^2/2=4$.

\begin{center}\adjustbox{max width=\linewidth}{$\displaystyle C=\frac{2\sqrt2}{\sqrt4}=\sqrt2$}\end{center}
\jegyzet{definíció}{}
\begin{eredmeny}

\begin{center}\adjustbox{max width=\linewidth}{$\displaystyle C=\sqrt2\qquad(\hat x=2\sqrt2,\ P=4)$}\end{center}
\end{eredmeny}
